\documentclass[11pt]{article}
\newcommand{\courseshort}{Data Structures (Make-up) · 010153103}
\newcommand{\lecturetitle}{L4: Stacks \& Queues}
\usepackage{lecture_style}

\begin{document}
\lectureheader{Lecture 4: Stacks, Queues \& Deques}{Pre-class brief}{Goodrich Ch.~6}

\begin{objbox}
\begin{itemize}[leftmargin=*]
  \item State the LIFO (stack) and FIFO (queue) principles and their core operations.
  \item Implement a stack with an array and with a linked list, both in $O(1)$.
  \item Implement a queue using a \emph{circular} array.
  \item Apply stacks/queues to real problems (matching, evaluation, BFS-like simulation).
\end{itemize}
\end{objbox}

\section*{1. Stack --- LIFO}

Operations: \texttt{push(e)}, \texttt{pop()}, \texttt{top()}, \texttt{size()}, \texttt{isEmpty()}.

\begin{lstlisting}
public class ArrayStack<E> {
    private E[] data;
    private int t = -1;                 // top index
    public ArrayStack(int cap) { data = (E[]) new Object[cap]; }
    public void push(E e) { data[++t] = e; }
    public E pop()        { E e = data[t]; data[t--] = null; return e; }
    public E top()        { return data[t]; }
    public boolean isEmpty() { return t < 0; }
}
\end{lstlisting}

A linked-list-based stack uses \texttt{addFirst}/\texttt{removeFirst} of an SLL, so all operations are $O(1)$ and the size is unbounded.

\section*{2. Queue --- FIFO}

Operations: \texttt{enqueue(e)}, \texttt{dequeue()}, \texttt{first()}, \texttt{size()}, \texttt{isEmpty()}.

\textbf{Circular array} avoids shifting on dequeue. Maintain \texttt{f} (index of front) and \texttt{sz} (current size); the back index is $(f+\text{sz}) \bmod N$.

\begin{lstlisting}
public void enqueue(E e) {
    if (sz == data.length) throw new IllegalStateException("full");
    int avail = (f + sz) % data.length;
    data[avail] = e;
    sz++;
}
public E dequeue() {
    E e = data[f];
    data[f] = null;
    f = (f + 1) % data.length;
    sz--;
    return e;
}
\end{lstlisting}

\section*{3. Deque (Double-Ended Queue)}

Allows insert/remove at \emph{both} ends in $O(1)$. A doubly linked list with sentinels is the canonical implementation.

\section*{4. Quick Reference}

\begin{center}
\begin{tabular}{lll}
\textbf{Structure} & \textbf{All ops} & \textbf{Notes} \\
\hline
Array stack         & $O(1)$ & fixed capacity \\
Linked stack        & $O(1)$ & dynamic \\
Circular array queue& $O(1)$ & fixed capacity \\
DLL deque           & $O(1)$ & dynamic, both ends \\
\end{tabular}
\end{center}

\begin{exbox}{Pre-Class Exercises (do BEFORE the lecture)}
\textbf{Exercise 4.1 --- Bracket matching.} Write a method
\begin{lstlisting}
boolean isBalanced(String s);
\end{lstlisting}
that returns \texttt{true} iff the brackets in \texttt{s} are balanced. Allowed pairs: \verb|()|, \verb|[]|, \verb|{}|. Use a stack. Examples: \verb|"({[]})"| $\to$ true, \verb|"([)]"| $\to$ false.

\textbf{Exercise 4.2 --- Reverse with a stack.} Use a stack to reverse the contents of an array in place. State the time and space complexity.

\textbf{Exercise 4.3 --- Postfix evaluation.} Evaluate \verb|"5 1 2 + 4 * + 3 -"| step by step using a stack. Show the stack contents after every token.

\textbf{Exercise 4.4 --- Trace circular queue.} Capacity 5, all initially empty. Show \texttt{f}, \texttt{sz}, and array contents after:
\verb|enq(A); enq(B); enq(C); deq(); deq(); enq(D); enq(E); enq(F);|.

\textbf{Exercise 4.5 --- Implement.} Implement a queue using \emph{two stacks}. \texttt{enqueue} should push to stack 1; \texttt{dequeue} should pop from stack 2 (refilling from stack 1 when empty). Show that the amortized cost of \texttt{dequeue} is $O(1)$.

\textbf{Exercise 4.6 --- When does a deque help?} Give an example task where a deque is strictly better than both a stack and a queue.
\end{exbox}

\begin{disbox}
\begin{itemize}[leftmargin=*]
  \item Why is the back index $(f+\text{sz}) \bmod N$? What goes wrong if you maintain a back pointer instead and just check \texttt{f == back}?
  \item Discuss Exercise 4.5 --- where exactly does the amortized argument break down if both stacks aren't used carefully?
  \item Real-world: how does the JVM call stack relate to recursion (preview of Lecture 5)?
\end{itemize}
\end{disbox}

\end{document}
